\documentclass[11pt]{article}
\usepackage[margin=1in]{geometry}
\usepackage{amsmath,amssymb}
\usepackage{microtype}
\usepackage{xcolor}
\usepackage[hidelinks]{hyperref}

\definecolor{deepblue}{RGB}{28,63,110}

\title{Dense Irrational Periodicity with a Single Rational Exception}
\author{}
\date{}

\begin{document}
\maketitle

There is an analytic Hamiltonian brake system in which a dense but
measure-zero set of irrational parameters is periodic, exactly one rational
parameter is periodic, and every other rational parameter is nonperiodic.

For \(0<u<1\), consider
\[
  H_u
  =\frac12\bigl(p_x^2+x^2\bigr)
   +\frac12\left(p_y^2+\Omega(u)^2y^2\right),
\]
where
\[
  \Omega(u)=2+u^2+\sqrt2\,u\left(u-\frac12\right).
\]
Start every orbit from the same brake state,
\[
  x(0)=y(0)=1,
  \qquad
  p_x(0)=p_y(0)=0.
\]
The resulting motion is
\[
  x(t)=\cos t,
  \qquad
  y(t)=\cos\!\bigl(\Omega(u)t\bigr).
\]
The full orbit is periodic---equivalently, it reaches a second brake---exactly
when
\[
  \Omega(u)\in\mathbb Q.
\]

\section*{Almost no rational parameters work}

Suppose \(u\in\mathbb Q\).  Then \(2+u^2\) is rational, whereas
\[
  \sqrt2\,u\left(u-\frac12\right)
\]
is irrational unless its rational coefficient vanishes.  Therefore
\[
  \Omega(u)\in\mathbb Q
  \quad\Longleftrightarrow\quad
  u\left(u-\frac12\right)=0.
\]
Since \(0<u<1\), the only rational periodic parameter is
\[
  \boxed{u=\frac12}.
\]
At that parameter,
\[
  \Omega\left(\frac12\right)=\frac94,
\]
so the orbit is genuinely periodic.  Every other rational parameter in the
interval produces quasiperiodic motion.

\section*{Nevertheless, irrational periodic parameters are dense}

Set
\[
  u_*=1-\frac1{\sqrt2}
      \approx0.292893218.
\]
This parameter is irrational, but direct substitution gives
\[
  \Omega(u_*)=2.
\]
It therefore generates the elementary \(1{:}2\) periodic orbit
\[
  x(t)=\cos t,
  \qquad
  y(t)=\cos(2t).
\]
Notably, this explicit irrational periodic parameter lies close to the
\(u\approx0.29\) region considered in the Pythagorean--Burrau investigation.

The complete periodic set is
\[
  \mathcal P=\Omega^{-1}(\mathbb Q).
\]
Because \(\Omega\) is nonconstant and continuous, the image of every open
parameter interval contains an interval and hence a rational number.  Thus
\(\mathcal P\) is dense in \((0,1)\).  On the other hand, every rational
frequency has at most two preimages under the quadratic function \(\Omega\).
Consequently, \(\mathcal P\) is countable and has Lebesgue measure zero.

Removing its sole rational member leaves
\[
  \mathcal P\setminus\left\{\frac12\right\},
\]
a dense, countable, measure-zero set consisting entirely of irrational periodic
parameters.

\begin{center}
  \color{deepblue}
  \fbox{\parbox{0.86\textwidth}{\centering
  Periodic irrational parameters are dense but measure zero; exactly one
  rational parameter is periodic; every other rational parameter is
  nonperiodic.}}
\end{center}

The example is deliberately engineered, but it shows exactly why measure,
density, and a generic dimension count cannot determine which rational
parameters lie on a periodic locus.  An exact structural or arithmetic bridge
is indispensable in the Pythagorean--Burrau problem.

\end{document}
